\section{Ablation and Verification Work}
We evaluate two complementary ablation levels. The policy-level study removes typed contracts (P0), keeps invalidation records without discharging obligations (P1), or discharges only an explicitly demanded obligation (P1D). P2 discharges every due obligation at its first eligible boundary. The mechanism-level study changes one validator behavior at a time while holding fixed one minimal counterexample and one valid control.

\begin{table}[t]
\caption{Policy-level ablation. ``Hist. W/T'' are the historical System W and TensorRules faults; ``Demand W/T'' are matched queried/unqueried demand cases. The table is regenerated from validated evidence and byte-compared in CI.}
\label{tab:policy-ablation}
\scriptsize
\centering
\begin{tabular}{@{}lrrrrr@{}}
\toprule
Variant & Prim. & Chall. & Hist. W & Hist. T & Demand W/T \\
\midrule
P0 no contracts & 12/32 & 1/10 & 0/5 & 0/8 & 0/2, 0/2 \\
P1 no discharge & 28/32 & 10/10 & 0/5 & 0/8 & 0/2, 0/2 \\
P1D demand only & 28/32 & 10/10 & 0/5 & 0/8 & 1/2, 1/2 \\
P2 selective & 32/32 & 10/10 & 5/5 & 8/8 & 2/2, 2/2 \\
\bottomrule
\end{tabular}

\end{table}

Removing typed contracts loses 20 primary and nine challenge detections relative to P2. Keeping invalidation without discharge recovers the challenge set but still loses four primary shapes and every targeted early rejection in both historical source-level studies. P1D retains the same historical behavior as P1, but rejects exactly the queried member of each matched demand pair; the unqueried member remains undischarged. P2 rejects both members because both obligations become due at their first eligible semantic boundary.

\begin{table}[t]
\caption{Eight isolated mechanism ablations. ``Full'' and ``Ablated'' report detection of the corresponding minimal counterexample; ``Control FP'' reports rejection of its valid control by the ablated validator. The table is evidence-generated and byte-compared in CI.}
\label{tab:mechanism-ablation}
\scriptsize
\centering
\begin{tabular}{@{}lp{0.43\columnwidth}ccc@{}}
\toprule
ID & Removed mechanism & Full & Ablated & Control FP \\
\midrule
M01 & Producer identity & 1/1 & 0/1 & 0/1 \\
M02 & Source identity & 1/1 & 0/1 & 0/1 \\
M03 & Selected order & 1/1 & 0/1 & 0/1 \\
M04 & Canonical owner & 1/1 & 0/1 & 0/1 \\
M05 & Route conflict & 1/1 & 0/1 & 0/1 \\
M06 & Fail-closed route & 1/1 & 0/1 & 0/1 \\
M07 & Capability contract & 1/1 & 0/1 & 0/1 \\
M08 & Repeated occurrence & 1/1 & 0/1 & 0/1 \\
\bottomrule
\end{tabular}

\end{table}

The full protocol detects all eight counterexamples. Removing each mechanism independently loses the matching detection, while neither the full nor ablated validators reject the eight controls. These are necessity witnesses for the implemented mechanisms, not estimates of natural fault frequency and not a proof that the mechanisms are jointly minimal for every compiler.

On 120 boundary controls, P2 executes 120 verifier calls and P3 executes 160, a 25\% reduction. This meets the frozen isolated-work threshold, but it is neither a pinned-machine result nor whole-compilation timing. We therefore report it only as verifier-work evidence and make no efficiency headline.
